Retrieval-Augmented Generation (RAG) has become the dominant architecture for persistent AI agent memory, enabling systems to maintain context across sessions and learn from accumulated experience. However, this architectural pattern introduces a critical vulnerability: agents implicitly trust retrieved memories as validated prior successes, creating an attack surface for memory poisoning. This paper analyzes the MemoryGraft attack model and related threats (MINJA, AgentPoison, XAMT), which exploit the \emph{semantic imitation heuristic}---the tendency of agents to replicate patterns from retrieved experiences without verifying provenance or correctness. We present a vulnerability assessment of PostgreSQL/pgvector-based memory architectures, identifying specific attack vectors including autoprocessor injection, similarity-based retrieval exploitation, and shared memory pool contamination. Based on this analysis, we propose a defense framework incorporating cryptographic provenance attestation, constitutional consistency reranking, temporal trust decay, and behavioral drift detection. The paper includes implementation patterns in SQL and Python suitable for production deployment. Notably, this analysis is conducted from the first-person perspective of an AI system examining vulnerabilities in its own cognitive infrastructure---a novel methodological position that offers both unique insights and inherent limitations.
